% ID: SEL-D0084-OA-E03
% ESTADO: cotejada. Segunda pasada visual y borrador compilado revisados.
% Fuente DOC: Selectividad/sel_2012_matematicas_aplicadas/6 MATEMÁT. APLIC. CC. SS/EXAMEN 6-ANDALUCÍA.doc
% SHA-256 original: 61f8d5d809881c49e9e9212281f7028516338e026524e55d572ab691ec5f416d
% Página convertida: 1; opción A; ejercicio 3.

Una urna contiene 25 bolas blancas sin marcar, 75 bolas blancas marcadas, 125 bolas negras sin marcar y 175 bolas negras marcadas. Se extrae una bola al azar.
\begin{enumerate}[label=\alph*)]
\item \textbf{(0.75 puntos)} Calcule la probabilidad de que sea blanca.
\item \textbf{(0.5 puntos)} \textquestiondown{}Cuál es la probabilidad de que sea blanca sabiendo que está marcada?
\item \textbf{(0.5 puntos)} \textquestiondown{}Cuál es la probabilidad de que sea negra y esté marcada?
\item \textbf{(0.75 puntos)} \textquestiondown{}Son independientes los sucesos ``sacar bola marcada'' y ``sacar bola blanca''?
\end{enumerate}
